\documentclass[11pt,a4paper]{article}
\usepackage[utf8]{inputenc}
\usepackage[T1]{fontenc}
\usepackage{lmodern}
\usepackage[margin=2.4cm]{geometry}
\usepackage{amsmath,amssymb,bm}
\usepackage{graphicx}
\usepackage{booktabs}
\usepackage[section]{placeins}
\usepackage{needspace}
\usepackage{xcolor}
\usepackage[font=small,labelfont=bf]{caption}
\usepackage[colorlinks=true,linkcolor=blue!50!black,citecolor=blue!50!black,urlcolor=blue!50!black]{hyperref}
\graphicspath{{figures/}}

\newcommand{\up}{{\uparrow}}
\newcommand{\dn}{{\downarrow}}
\newcommand{\I}{\mathrm{i}}
\newcommand{\Tr}{\operatorname{Tr}}
\newcommand{\todo}[1]{\textcolor{red}{[#1]}}
% Figures from figures/ (PNG, written by the plotting scripts); a frame marks a missing one.
\newcommand{\fig}[1]{\IfFileExists{figures/#1.png}{\includegraphics[width=\textwidth]{#1.png}}%
  {\fbox{\parbox[c][4cm][c]{0.9\textwidth}{\centering missing figure \texttt{\detokenize{#1}}}}}}

\title{The sign of the mass self-energy in the Born analysis of\\
``Anomalous quantum Hall effect induced by disorder in topological insulators''}
\author{Arnaud Demion}
\date{October 2026}

\begin{document}
\maketitle

\begin{abstract}
The self-consistent Born approximation (SCBA) used in Ref.~\cite{g2g1}, and in Chapter~2 of the
thesis~\cite{thesis}, renormalises the masses of the two spin blocks with the wrong sign. The
error is in the definition of the renormalised (barred) parameters,
$\mathrm{Re}\,\Sigma \equiv -\mathrm{diag}[(\bar m - m)\tau_z + (\bar\mu - \mu)\tau_0,\dots]$, which holds for
the chemical potential but not for the mass. The program that computed the Born phase diagrams
(Fig.~4 of Ref.~\cite{g2g1}) implements these equations as written. With the correct sign, the
disorder fluctuations push both masses towards the topological side, as in the topological
Anderson insulator that the paper follows. Three checks confirm it, none of which uses the
derivation:
a full matrix SCBA, the exact disorder-averaged Green's function, and spin-resolved transport. The
numerical results of the paper (Figs.~1--3) are unaffected, and so is its main conclusion: a
$G=1$ Chern-insulator phase, carried by the spin-down edge state, for both couplings. What changes
is the Born interpretation. In the same-$g$ (``ferromagnetic'') case, the spin-down block, which
carries the $G=1$ plateau, never changes sign, whereas the published lines invert it. The
spin-up channel is lost because the Fermi level enters the spin-up band, shifted by the mean
exchange field. At stronger disorder the spin-up block does become trivial, as the paper says,
but through a resonance that no Born approximation describes. In the opposite-$g$ (``antiferromagnetic'') case, the spin-up mass changes
sign at $J_I\simeq 0.40$, not $0.30$. The apparent agreement of the published inversion line
with the antiferromagnetic conductance map is a finite-width effect of the 64-site ribbons.
\end{abstract}

\tableofcontents

%======================================================================
\section{Summary}

\begin{itemize}
\item \textbf{The derivation} (Sec.~\ref{s:derivation}). Every step of the SCBA, from the
  Hamiltonian to the program; the error is in \hyperref[st:identify]{Step~7}.
\item \textbf{The error.} The mass equations of the SCBA carry the wrong sign
  (Sec.~\ref{s:error}). The chemical-potential equations are right.
\item \textbf{Where it comes from.} The derivation, as printed in Ref.~\cite{g2g1} and in
  the thesis~\cite{thesis}: the code \texttt{self\_born.f90} implements the printed equations
  faithfully. A replica of the program reproduces the published Born data (runs 00--07) to
  $10^{-14}$.
\item \textbf{Independent evidence} (Sec.~\ref{s:checks}):
  \begin{enumerate}
  \item a matrix SCBA computed from the Hamiltonian, without the reduction;
  \item the exact renormalised masses, from the disorder-averaged Green's function
    (Fig.~\ref{f:masses});
  \item the spin-resolved conductance of the same ribbons as the paper.
  \end{enumerate}
  All three follow the corrected sign. At weak disorder the corrected SCBA agrees with the exact
  masses to $10^{-3}$, while the published one moves them in the opposite direction.
\item \textbf{Impact on Fig.~4.} With the published maps (Fig.~\ref{f:phase},
  Sec.~\ref{s:impact}):
  \begin{itemize}
  \item same-$g$: the corrected lines explain $99.9\,\%$ of the $G=1$ plateau, the published
    ones $20\,\%$;
  \item opposite-$g$: the published spin-up inversion line sits at the left edge of the $G=1$
    region, the corrected one inside it.
  \end{itemize}
\item \textbf{Why the published lines looked right} (Sec.~\ref{s:width}). The ribbons are 64
  sites wide. Close to the inversion, the spin-up edge states decay over $1.5/|\bar m_\up|$
  sites, so they overlap across the ribbon and the spin-up channel is lost before the mass
  changes sign. Wider ribbons move the transition towards the exact inversion. At $W=256$ the
  spin-up block still carries $T_\up=0.99$ where the published masses make it trivial. The
  wrong sign moves the Born inversion by about the same amount in the same direction: a
  coincidence.
\item \textbf{Origin and stability of $G=1$} (Sec.~\ref{s:origin}). $G=1$ is a Chern
  insulator: the spin-down block keeps its topological gap, while the spin-up states at the
  Fermi level no longer conduct. Each edge then carries a single chiral mode, with nothing to
  backscatter into. The same-$g$ plateau does not depend on the ribbon width, and it survives
  tilted moments ($M_z=0.8$), whereas $G=2$ does not.
\end{itemize}

%======================================================================
\section{The model and the published equations}\label{s:model}

The two spin blocks of the BHZ model are, on the square lattice (lattice constant 1),
\begin{equation}
h_\sigma(\bm k) = -D_{\bm k}\,\tau_0 + (m - B_{\bm k})\,\tau_z + A\,(\text{terms odd in }\bm k),
\qquad B_{\bm k} = B(4-2\cos k_x-2\cos k_y),
\label{e:h}
\end{equation}
and similarly for $D_{\bm k}$. The parameters are $A=1$, $B=-0.376406$, $D=-0.280933$ and
$m=-0.137$, so the clean system is topological, $m/B>0$. The Pauli matrices $\tau$ act on the
electron (E) and hole (H) orbitals.

The magnetic impurities occupy a fraction $n_I$ of the sites. They carry classical moments
$\bm S_i$ with $\langle S_z\rangle = M_z$ and couple through
\begin{equation}
V_\alpha = J_I \sum_{i\in I} \bm S_i\cdot\bm\sigma\,\tau_\alpha ,\qquad
\tau_\alpha = \tau_0\ \text{(same-$g$, ``F'')},\quad \tau_\alpha=\tau_z\ \text{(opposite-$g$, ``AF'')}.
\label{e:V}
\end{equation}
This note treats the fully polarised case $M_z=1$ of the paper's Figs.~1, 2 and~4. Then
$V_\alpha = J_I\sigma_z\tau_\alpha$ on the impurity sites, $s_z$ is conserved, and the two spin blocks
are independent disordered systems. The mean field $\langle V_\alpha\rangle = n_I J_I M_z\,\sigma_z\tau_\alpha$
shifts, block by block:
\begin{itemize}
\item the chemical potential, $\mu_\sigma = \mu - \sigma n_I J_I M_z$, for same-$g$;
\item the mass, $m_\sigma = m + \sigma n_I J_I M_z$, for opposite-$g$.
\end{itemize}

Ref.~\cite{g2g1} (Secs.~IV\,A ``Born approximation'' and IV\,B ``Phase diagram'') and the
thesis~\cite{thesis} (Sec.\ ``Discussion : l'approximation de Born'') write the mean-field Green's function and the SCBA self-energy as
\begin{equation}
g_\alpha(E,\bm k) = \left[E+\mu-H_0(\bm k)-\langle V_\alpha\rangle\right]^{-1},\qquad
\Sigma(E) = \frac{n_I J_I^2}{N^2}\sum_{\bm k}\left[v_z\, g_z + v_\perp g_\perp\right],
\label{e:sigma}
\end{equation}
where $g$ is replaced by the full averaged Green's function in the self-consistent version. They
then \emph{define} the renormalised parameters by
\begin{equation}
\mathrm{Re}\,\Sigma \equiv -\mathrm{diag}\big[(\bar m_\up-m_\up)\tau_z+(\bar\mu_\up-\mu_\up)\tau_0,\;
(\bar m_\dn-m_\dn)\tau_z-(\bar\mu_\dn-\mu_\dn)\tau_0\big]
\qquad\text{(as printed)},
\label{e:def}
\end{equation}
and obtain
\begin{align}
\bar\mu_\sigma &= \mu_\sigma - \frac{n_I J_I^2}{N^2}\sum_{\bm k}\Big[\frac{v_z}{\Delta_\sigma}(E_F+\bar\mu_\sigma+D_{\bm k})
 + \frac{v_\perp}{\Delta_{\bar\sigma}}(E_F+\bar\mu_{\bar\sigma}+D_{\bm k})\Big],\label{e:mu}\\
\bar m_\sigma &= m_\sigma - \frac{n_I J_I^2}{N^2}\sum_{\bm k}\Big[\frac{v_z}{\Delta_\sigma}(\bar m_\sigma-B_{\bm k})
 + \frac{v_\perp}{\Delta_{\bar\sigma}}(\bar m_{\bar\sigma}-B_{\bm k})\Big],\label{e:m}
\end{align}
with $\Delta_\sigma = (E_F+\bar\mu_\sigma+D_{\bm k})^2-(\bar m_\sigma-B_{\bm k})^2-\sin^2k_x-\sin^2k_y$. The
gap of block $\sigma$ is centred on $-\bar\mu_\sigma$, with half-width $|\bar m_\sigma|$, and the block
is topological when $\bar m_\sigma/B>0$. Fig.~4 of Ref.~\cite{g2g1} draws the band edges
$E_F=-\bar\mu_\sigma\pm\bar m_\sigma$ and the lines $\bar m_\sigma=0$.\footnote{Eq.~(\ref{e:def}) as
printed also has opposite signs for the $\tau_0$ terms of the two blocks. Eqs.~(\ref{e:mu}) use
the same sign for both, which is the right one. The thesis writes $\mu_\up$ instead of $m_\up$
on the right-hand side of the spin-up mass equation; the program uses $m_\up$.}

The program \texttt{self\_born.f90}, which produced Fig.~4, iterates exactly
Eqs.~(\ref{e:mu}) and~(\ref{e:m}). In it, \texttt{norm} is $n_I J_I^2/N^2$, \texttt{kc} is
$B_{\bm k}/B$, and \texttt{vzu = vz/detu} with \texttt{detu} $=\Delta_\up$:
\begin{verbatim}
  154  vz = (1-conc**2)*Mz**2 + (1-Mz)**2/3.0
  192  deltaMu = -norm * sum( vzu*(Mu - B*kc) + vperpd*(Md - B*kc) )
  200  Mu_temp = Mu0 + deltaMu
\end{verbatim}
A replica written for this note reproduces the published Born data of runs 00--07 to
$10^{-14}$. The code is therefore a faithful implementation of the printed equations.

%======================================================================
\section{The self-consistent Born approximation, step by step}\label{s:derivation}

This section rederives the SCBA of Ref.~\cite{g2g1} from the Hamiltonian, with every
intermediate step, in the paper's notation. It departs from the printed equations in two places:
the identification of the renormalised mass (\hyperref[st:identify]{Step~7}) and the weight $v_z$
(\hyperref[st:disorder]{Step~2}). Both are flagged where they occur. Spins are labelled $\sigma=+1$ ($\up$)
and $\sigma=-1$ ($\dn$), and $\bar\sigma=-\sigma$.

\subsection{Step 1: the clean model in \texorpdfstring{$\bm k$}{k} space}\label{st:clean}

Each site carries four orbitals, $(E\up, H\up, E\dn, H\dn)$. The Pauli matrices $\bm\sigma$ act on
the spin and $\bm\tau$ on the E/H orbitals. After a Fourier transform on the square lattice, the
clean Hamiltonian is block diagonal in spin, $H_0(\bm k)=\mathrm{diag}[h_\up(\bm k),h_\dn(\bm k)]$, with
\begin{equation}
h_\sigma(\bm k) = -D_{\bm k}\,\tau_0 + (m-B_{\bm k})\,\tau_z + A\,(\sigma\sin k_x\,\tau_x - \sin k_y\,\tau_y),
\qquad B_{\bm k}=B\,(4-2\cos k_x-2\cos k_y),
\label{e:hk}
\end{equation}
and the same for $D_{\bm k}$. Two properties of the $\tau_x$, $\tau_y$ terms are used below:
\begin{itemize}
\item they are odd in $\bm k$;
\item their squares add up to $A^2(\sin^2k_x+\sin^2k_y)$.
\end{itemize}
The chemical potential enters as $H_0-\mu$, so Green's functions are $[E+\mu-H]^{-1}$. In this
form the mass appears with a plus sign ($+m\,\tau_z$) and the chemical potential with a minus
sign ($-\mu\,\tau_0$). This difference is what \hyperref[st:identify]{Step~7} is about.

\subsection{Step 2: the disorder and its first two moments}\label{st:disorder}

Each site $\bm r$ is occupied by an impurity with probability $n_I$, independently of the other
sites: $\eta_{\bm r}\in\{0,1\}$, $\langle\eta_{\bm r}\rangle=n_I$. An impurity carries a classical unit
moment $\bm S_{\bm r}$, independent of $\eta$, with
\begin{itemize}
\item a uniform azimuth, so $\langle S_x\rangle=\langle S_y\rangle=0$ and $\langle S_aS_b\rangle=0$ for $a\ne b$;
\item $S_z$ uniform in $[2M_z-1,1]$ (the polarised moments of the thesis), so $\langle S_z\rangle=M_z$
  and $\langle S_z^2\rangle = M_z^2+(1-M_z)^2/3$;
\item $\langle S_x^2\rangle=\langle S_y^2\rangle=(1-\langle S_z^2\rangle)/2$, since $|\bm S|=1$.
\end{itemize}
The potential on site $\bm r$ is
\begin{equation}
U_{\bm r} = J_I\,\eta_{\bm r}\,(\bm S_{\bm r}\cdot\bm\sigma)\,\tau_\alpha = J_I\sum_{a=x,y,z}X^{\bm r}_a\,\sigma_a\tau_\alpha,
\qquad X_a=\eta\,S_a .
\end{equation}
Its mean is $\langle X_a\rangle = n_I\langle S_a\rangle$, so
\begin{equation}
\langle U\rangle = n_I J_I M_z\,\sigma_z\tau_\alpha .
\end{equation}
Its covariances follow from $\eta^2=\eta$:
$C_{ab}=\langle X_aX_b\rangle-\langle X_a\rangle\langle X_b\rangle = n_I\langle S_aS_b\rangle-n_I^2\langle S_a\rangle\langle S_b\rangle$. Only the
diagonal ones survive:
\begin{equation}
C_{zz} = n_I\big(\langle S_z^2\rangle-n_IM_z^2\big)\equiv n_I\,v_z,\qquad
C_{xx}=C_{yy}=\tfrac{n_I}{2}\big(1-\langle S_z^2\rangle\big)\equiv \tfrac{n_I}{2}\,v_\perp ,
\end{equation}
\begin{equation}
v_z=(1-n_I)M_z^2+\tfrac13(1-M_z)^2,\qquad v_\perp=\tfrac23(1-M_z)(1+2M_z).
\label{e:weights}
\end{equation}
Ref.~\cite{g2g1} writes $v_z=\langle S_z^2\rangle-n_I^2M_z^2$: the factor $n_I$ is taken out twice
(Sec.~\ref{s:vz}). $v_\perp$ agrees.

\subsection{Step 3: averaged Green's function, mean field and self-energy}\label{st:dyson}

For one configuration, $G=[z+\mu-H_0-V]^{-1}$, with $z=E+\I0^+$ and $V=\sum_{\bm r}U_{\bm r}$. The disorder
average restores translation invariance, so $\langle G\rangle$ is diagonal in $\bm k$. The self-energy is
defined by the Dyson equation:
\begin{equation}
\langle G\rangle(\bm k) = \big[z+\mu-H_0(\bm k)-\langle V\rangle-\Sigma(\bm k,z)\big]^{-1}.
\label{e:dyson}
\end{equation}
The mean field $\langle V\rangle=\langle U\rangle$ is uniform. As in the paper, it goes into the unperturbed
Green's function, $g(\bm k)=[z+\mu-H_0(\bm k)-\langle V\rangle]^{-1}$. Block by block:
\begin{itemize}
\item same-$g$ ($\tau_\alpha=\tau_0$): $\langle V\rangle_\sigma=\sigma n_IJ_IM_z\,\tau_0$, so $\mu_\sigma=\mu-\sigma n_IJ_IM_z$ and
  $m_\sigma=m$;
\item opposite-$g$ ($\tau_\alpha=\tau_z$): $\langle V\rangle_\sigma=\sigma n_IJ_IM_z\,\tau_z$, so $\mu_\sigma=\mu$ and
  $m_\sigma=m+\sigma n_IJ_IM_z$.
\end{itemize}
Then $g_\sigma(\bm k)=[z+\mu_\sigma-h_\sigma(\bm k;m_\sigma)]^{-1}$, where $h_\sigma(\bm k;m_\sigma)$ is
Eq.~(\ref{e:hk}) with $m$ replaced by $m_\sigma$.

\subsection{Step 4: second order, the Born approximation}\label{st:born}

Write $\delta U_{\bm r}=U_{\bm r}-\langle U\rangle$ and expand: $G=g+g\,\delta V g+g\,\delta Vg\,\delta Vg+\dots$ Since
$\langle\delta V\rangle=0$,
\begin{equation}
\langle G\rangle = g + g\,\langle\delta V g\,\delta V\rangle\,g + O(\delta V^3)
\quad\Longrightarrow\quad \Sigma=\langle\delta V g\,\delta V\rangle + O(\delta V^3).
\end{equation}
On the sites, $(\delta Vg\,\delta V)_{\bm r\bm r'}=\delta U_{\bm r}\,g_{\bm r\bm r'}\,\delta U_{\bm r'}$. Different sites are independent
and $\delta U$ has zero mean, so the average vanishes unless $\bm r=\bm r'$:
\begin{equation}
\Sigma_{\bm r\bm r'}=\delta_{\bm r\bm r'}\,\langle\delta U\,g_{\mathrm{loc}}\,\delta U\rangle,\qquad
g_{\mathrm{loc}}=g_{\bm r\bm r}=\frac1{N_k}\sum_{\bm k}g(\bm k).
\end{equation}
The self-energy is local, hence independent of $\bm k$. $N_k$ is the number of $\bm k$ points
($N^2$ in the paper). The self-consistent version replaces $g$ by $\langle G\rangle$ itself, which resums
all the non-crossing diagrams:
\begin{equation}
\Sigma=\langle\delta U\,\langle G\rangle_{\mathrm{loc}}\,\delta U\rangle,\qquad
\langle G\rangle_{\mathrm{loc}}=\frac1{N_k}\sum_{\bm k}\big[z+\mu-H_0(\bm k)-\langle V\rangle-\Sigma\big]^{-1}.
\label{e:scba}
\end{equation}

\subsection{Step 5: the vertex algebra}\label{st:vertex}

With the covariances of \hyperref[st:disorder]{Step~2},
\begin{equation}
\Sigma = J_I^2\sum_{ab}C_{ab}\,(\sigma_a\tau_\alpha)\langle G\rangle_{\mathrm{loc}}(\sigma_b\tau_\alpha)
= n_IJ_I^2\Big[v_z\,(\sigma_z\tau_\alpha)\langle G\rangle_{\mathrm{loc}}(\sigma_z\tau_\alpha)
+\frac{v_\perp}{2}\sum_{a=x,y}(\sigma_a\tau_\alpha)\langle G\rangle_{\mathrm{loc}}(\sigma_a\tau_\alpha)\Big].
\end{equation}
Take $\langle G\rangle$ block diagonal in spin; Eq.~(\ref{e:scba}) then keeps it so. For a
block-diagonal matrix,
\begin{equation}
\sigma_z\,\mathrm{diag}(P,Q)\,\sigma_z=\mathrm{diag}(P,Q),\qquad
\sigma_x\,\mathrm{diag}(P,Q)\,\sigma_x=\sigma_y\,\mathrm{diag}(P,Q)\,\sigma_y=\mathrm{diag}(Q,P).
\end{equation}
Longitudinal fluctuations keep the spin and transverse ones exchange the blocks:
\begin{equation}
\Sigma_\sigma = n_IJ_I^2\big[v_z\,\tau_\alpha\langle G_\sigma\rangle_{\mathrm{loc}}\tau_\alpha+v_\perp\,\tau_\alpha\langle G_{\bar\sigma}\rangle_{\mathrm{loc}}\tau_\alpha\big].
\end{equation}
\hyperref[st:local]{Step~6} shows that $\langle G_\sigma\rangle_{\mathrm{loc}}$ has only $\tau_0$ and $\tau_z$ components. These
commute with $\tau_z$, so $\tau_\alpha$ drops out for both couplings:
\begin{equation}
\Sigma_\sigma = n_IJ_I^2\big[v_z\,\langle G_\sigma\rangle_{\mathrm{loc}}+v_\perp\,\langle G_{\bar\sigma}\rangle_{\mathrm{loc}}\big].
\label{e:sigmablock}
\end{equation}
At this order the two couplings differ only through the mean field of \hyperref[st:dyson]{Step~3}.

\subsection{Step 6: the local Green's function of a block}\label{st:local}

By Eq.~(\ref{e:sigmablock}), $\Sigma_\sigma=\Sigma^0_\sigma\tau_0+\Sigma^z_\sigma\tau_z$. The averaged block is then
\begin{equation}
\langle G_\sigma\rangle(\bm k)=\big[a_\sigma(\bm k)-\bm d_\sigma(\bm k)\cdot\bm\tau\big]^{-1},\qquad
a_\sigma=z+\mu_\sigma+D_{\bm k}-\Sigma^0_\sigma,\quad
d^z_\sigma=m_\sigma-B_{\bm k}+\Sigma^z_\sigma,
\end{equation}
with $d^x_\sigma=\sigma A\sin k_x$ and $d^y_\sigma=-A\sin k_y$. Since $(\bm d\cdot\bm\tau)^2=\bm d\cdot\bm d$,
\begin{equation}
\big[a-\bm d\cdot\bm\tau\big]^{-1}=\frac{a+\bm d\cdot\bm\tau}{a^2-\bm d\cdot\bm d},\qquad
\Delta_\sigma(\bm k)=a_\sigma^2-\bm d_\sigma\cdot\bm d_\sigma .
\end{equation}
The sum over $\bm k$ removes the odd components $d^x$, $d^y$:
\begin{equation}
\langle G_\sigma\rangle_{\mathrm{loc}}=\frac1{N_k}\sum_{\bm k}\frac{a_\sigma(\bm k)\,\tau_0+d^z_\sigma(\bm k)\,\tau_z}{\Delta_\sigma(\bm k)} .
\label{e:gloc}
\end{equation}

\subsection{Step 7: the renormalised parameters}\label{st:identify}

The barred parameters are defined so that the averaged block is a clean block with new
parameters, $a_\sigma=z+\bar\mu_\sigma+D_{\bm k}$ and $d^z_\sigma=\bar m_\sigma-B_{\bm k}$:
\begin{equation}
\bar\mu_\sigma=\mu_\sigma-\Sigma^0_\sigma,\qquad \bar m_\sigma=m_\sigma+\Sigma^z_\sigma,
\qquad\text{i.e.}\qquad
\mathrm{Re}\,\Sigma_\sigma=(\bar m_\sigma-m_\sigma)\,\tau_z-(\bar\mu_\sigma-\mu_\sigma)\,\tau_0 .
\end{equation}
\textbf{This is where the printed equations go wrong.} Eq.~(\ref{e:def}) puts a minus sign in
front of both terms, i.e.\ $\bar m_\sigma=m_\sigma-\Sigma^z_\sigma$. The chemical potential enters
$z+\mu-h$ with a plus sign, and the mass enters through $-h$, i.e.\ with a minus sign. Their
self-energy corrections therefore have opposite signs. With the barred parameters,
$\Delta_\sigma=(z+\bar\mu_\sigma+D_{\bm k})^2-(\bar m_\sigma-B_{\bm k})^2-A^2(\sin^2k_x+\sin^2k_y)$, the
paper's determinant.

\subsection{Step 8: the four self-consistent equations}\label{st:equations}

Inserting Eq.~(\ref{e:gloc}) into Eq.~(\ref{e:sigmablock}), component by component:
\begin{align}
\Sigma^0_\sigma&=\frac{n_IJ_I^2}{N_k}\sum_{\bm k}\Big[\frac{v_z}{\Delta_\sigma}(z+\bar\mu_\sigma+D_{\bm k})
 +\frac{v_\perp}{\Delta_{\bar\sigma}}(z+\bar\mu_{\bar\sigma}+D_{\bm k})\Big],\\
\Sigma^z_\sigma&=\frac{n_IJ_I^2}{N_k}\sum_{\bm k}\Big[\frac{v_z}{\Delta_\sigma}(\bar m_\sigma-B_{\bm k})
 +\frac{v_\perp}{\Delta_{\bar\sigma}}(\bar m_{\bar\sigma}-B_{\bm k})\Big].
\end{align}
Then $\bar\mu_\sigma=\mu_\sigma-\Sigma^0_\sigma$, which is Eq.~(\ref{e:mu}) as printed, and
$\bar m_\sigma=m_\sigma+\Sigma^z_\sigma$, which is Eq.~(\ref{e:mok}) and differs from the printed Eq.~(\ref{e:m}) by the
sign. These are four equations for the four complex unknowns $(\bar\mu_\up,\bar\mu_\dn,\bar m_\up,\bar m_\dn)$
at each $(J_I,E_F)$, with $z=E_F+\I\eta$. For $M_z=1$, $v_\perp=0$ and the two blocks decouple.

\subsection{Step 9: band edges and topology, the lines of Fig.~4}\label{st:lines}

The renormalised block has its gap at $\bm k=0$, where $D_{\bm k}=B_{\bm k}=\sin k=0$: its bands start at
$E=-\bar\mu_\sigma\pm\bar m_\sigma$ (real parts). Hence:
\begin{itemize}
\item the band edges are $E_F=-\bar\mu_\sigma\pm\bar m_\sigma$. Since the parameters depend on $E_F$
  through $z$, each line of Fig.~4 is the zero contour of $E_F+\bar\mu_\sigma\mp\bar m_\sigma$ on the
  $(J_I,E_F)$ grid;
\item $E_F$ lies in the gap of block $\sigma$ when $|E_F+\bar\mu_\sigma|<|\bar m_\sigma|$;
\item block $\sigma$ is topological when $\bar m_\sigma/B>0$, a band inversion at $\Gamma$ (on the lattice,
  for $0<\bar m_\sigma/B<4$). The line $\bar m_\sigma=0$ is its topological transition.
\end{itemize}

\subsection{Step 10: lowest order, the logarithm and the sign of the mass correction}\label{st:log}

Put the mean-field parameters on the right-hand side of \hyperref[st:equations]{Step~8} and go to the
continuum, $N_k^{-1}\sum_{\bm k}\to\int_{|\bm k|<k_*}d^2k/(2\pi)^2$ with $k_*=\pi$:
\begin{equation}
\Sigma^z_\sigma\simeq n_IJ_I^2v_z\int\!\frac{d^2k}{(2\pi)^2}\,\frac{m_\sigma-Bk^2}{\Delta_\sigma(\bm k)},\qquad
\Delta_\sigma(\bm k)=(E_F+\mu_\sigma+Dk^2)^2-(m_\sigma-Bk^2)^2-A^2k^2 .
\end{equation}
At large $k$, $\Delta_\sigma\simeq-(B^2-D^2)k^4$ and the integrand tends to $B/[(B^2-D^2)k^2]$. Its integral
is logarithmic:
\begin{equation}
\int_{k_0}^{k_*}\frac{k\,dk}{2\pi}\,\frac{B}{(B^2-D^2)k^2}=\frac{B}{2\pi(B^2-D^2)}\,\ln\frac{k_*}{k_0}.
\end{equation}
It is cut off at small $k$, where $(B^2-D^2)k_0^4\simeq|\Delta_\sigma(E_F,0)|$. Hence
\begin{equation}
\mathrm{Re}\,\Sigma^z_\sigma\simeq\frac{n_IJ_I^2v_z}{8\pi}\,\frac{B}{B^2-D^2}\,
\ln\Big|\frac{(B^2-D^2)k_*^4}{\Delta_\sigma(E_F,0)}\Big| ,
\end{equation}
and in the same way $\mathrm{Re}\,\Sigma^0_\sigma\simeq-\frac{n_IJ_I^2v_z}{8\pi}\frac{D}{B^2-D^2}\ln|\dots|$. This is
Eq.~(\ref{e:reS}) of Sec.~\ref{s:error}, the paper's lowest-order result.

\paragraph{Sign.} $B^2-D^2>0$, and the logarithm is large and positive when $E_F$ is in the gap, so
$\mathrm{Re}\,\Sigma^z$ has the sign of $B$. With \hyperref[st:identify]{Step~7}, $\bar m_\sigma-m_\sigma=\mathrm{Re}\,\Sigma^z_\sigma$:
the mass moves away from zero on the topological side. With the printed sign, it moves towards
the trivial side.

\paragraph{Numbers.} Take same-$g$, $n_I=0.2$, $J_I=0.3$, $E_F=0$ and spin up, so
$\mu_\up=-0.06$. Then $B/(B^2-D^2)=-6.0$, the logarithm is $6.0$ and $\mathrm{Re}\,\Sigma^z\simeq-0.021$.
\begin{itemize}
\item Corrected sign: $\bar m_\up\simeq-0.158$, against $-0.149$ for the full corrected SCBA and
  $-0.153$ for the exact mass (Table~\ref{t:masses}).
\item Printed sign: $\bar m_\up\simeq-0.112$, against $-0.123$ for the published SCBA.
\end{itemize}

\subsection{Step 11: the numerical solution in the program}\label{st:program}

\texttt{self\_born.f90} solves \hyperref[st:equations]{Step~8} at each point of a $64\times64$ grid of
$(J_I,E_F)$:
\begin{itemize}
\item the $\bm k$ sum runs over a $256\times64$ grid of the Brillouin zone, with
  \texttt{norm} $=n_IJ_I^2/(256\cdot64)$, \texttt{kc} $=4-2\cos k_x-2\cos k_y$ and
  \texttt{ks} $=\sin^2k_x+\sin^2k_y$;
\item $z=E_F+\I\eta$ with $\eta=10^{-3}$;
\item the unknowns are \texttt{Mu}, \texttt{Md} ($\bar m_\up$, $\bar m_\dn$) and \texttt{Muu},
  \texttt{Mud} ($\bar\mu_\up$, $\bar\mu_\dn$). They start from the mean field,
  \texttt{Mu0 = mass + anti*J\_I*conc*Mz} and \texttt{Muu0 = -(1-anti)*J\_I*conc*Mz}, with
  \texttt{anti} $=1$ for opposite-$g$;
\item the iteration is a damped fixed point, $X\leftarrow pF(X)+(1-p)X$: first with $p=0.15$
  down to a tolerance $10^{-4}$ (at most 500 iterations), then with $p=0.8$ down to $10^{-6}$;
\item the program writes the real parts of $\bar m_\sigma$ and $\bar\mu_\sigma$.
\end{itemize}
The replica in \texttt{caroli} (\texttt{ScbaAnalytic}) reproduces it to $10^{-14}$.
\texttt{ScbaAnalytic::point\_corrected} changes only the sign of \hyperref[st:identify]{Step~7} and the
weight of \hyperref[st:disorder]{Step~2}.

%======================================================================
\section{The error}\label{s:error}

\subsection{Derivation}

Sec.~\ref{s:derivation} gives every step; here is the one that matters (\hyperref[st:identify]{Step~7}).
The averaged Green's function is
$\langle G\rangle = [E+\mu-H_0-\langle V\rangle-\Sigma]^{-1}$. The renormalised parameters
are defined by writing it as the Green's function of a clean block with barred parameters:
\begin{equation}
E+\mu_\sigma-h_\sigma(\bm k; m_\sigma)-\Sigma_\sigma = E+\bar\mu_\sigma - h_\sigma(\bm k;\bar m_\sigma).
\end{equation}
The mass enters $h_\sigma$ as $+m\,\tau_z$ (Eq.~\ref{e:h}), the chemical potential as $-\mu\,\tau_0$.
Identifying the $\tau_0$ and $\tau_z$ components gives
\begin{equation}
\boxed{\;\mathrm{Re}\,\Sigma_\sigma = (\bar m_\sigma-m_\sigma)\,\tau_z-(\bar\mu_\sigma-\mu_\sigma)\,\tau_0\;}
\label{e:defok}
\end{equation}
The minus sign of Eq.~(\ref{e:def}) is right for $\mu$ and wrong for $m$. The $\tau_z$ component of
$g_\sigma$ is $(\bar m_\sigma-B_{\bm k})/\Delta_\sigma$ (Ref.~\cite{g2g1}, Eq.~for $g_z$). The mass equation
therefore reads
\begin{equation}
\boxed{\;\bar m_\sigma = m_\sigma + \frac{n_I J_I^2}{N^2}\sum_{\bm k}\Big[\frac{v_z}{\Delta_\sigma}(\bar m_\sigma-B_{\bm k})
 + \frac{v_\perp}{\Delta_{\bar\sigma}}(\bar m_{\bar\sigma}-B_{\bm k})\Big]\;}
\label{e:mok}
\end{equation}
Eq.~(\ref{e:mu}) for $\bar\mu_\sigma$ is unchanged.

\subsection{One-line check}

Treat a uniform, non-random term $\delta\,\tau_z$ as a self-energy: $\Sigma=\delta\,\tau_z$ exactly, and
it shifts the mass to $m+\delta$. Eq.~(\ref{e:def}) gives $\bar m = m-\delta$.

\subsection{It is not a sign convention}

Some authors write $\langle G\rangle=[g^{-1}+\Sigma]^{-1}$ with $\Sigma=-\langle VgV\rangle$. A change of convention
flips the $\tau_0$ and $\tau_z$ components of $\Sigma$ together. Eqs.~(\ref{e:mu}) and~(\ref{e:m}) treat them
differently: both carry a minus sign in front of the sum, while $\mu$ enters the Hamiltonian
with a minus sign and $m$ with a plus sign. No convention makes both equations right; the
checks of Sec.~\ref{s:checks} show which one is wrong.

\subsection{Consistency with the perturbative result and with the topological Anderson insulator}

At lowest order, Ref.~\cite{g2g1} obtains
\begin{equation}
\mathrm{Re}\,\Sigma = \frac{n_I J_I^2}{8\pi}\,\frac{B\tau_z-D\tau_0}{B^2-D^2}
\Big[v_z\log\Big|\frac{(B^2-D^2)k_*^4}{\Delta(E_F,0)}\Big| + \dots\Big],
\label{e:reS}
\end{equation}
with a large positive logarithm when $E_F$ is in the gap. With Eq.~(\ref{e:defok}), the mass
correction $\bar m-m=\mathrm{Re}\,\Sigma_{\tau_z}$ has the sign of $B$. It moves $\bar m$ away from
zero on the topological side, $\bar m/B>0$. This is the renormalisation that explains the
topological Anderson insulator~\cite{li2009,groth2009,jiang2009}, the mechanism the paper says it
follows. With Eq.~(\ref{e:def}), the published equations push the masses towards the trivial side
instead.

Physically, at second order the two couplings act alike on each block: $\tau_0$ disorder (same-$g$)
and $\tau_z$ disorder (opposite-$g$) give the same $\langle VgV\rangle$ once the terms odd in $\bm k$ are
summed out, since $\tau_z(a+b\tau_z)\tau_z=a+b\tau_z$. Both are the scalar disorder of the topological
Anderson insulator.

\subsection{A second, smaller correction}\label{s:vz}

The longitudinal weight is printed as $v_z = \langle S_z^2\rangle - n_I^2M_z^2 = (1-n_I^2)M_z^2+(1-M_z)^2/3$.
The prefactor $n_I$ of Eq.~(\ref{e:sigma}) is already factored out of the second cumulant of
the Bernoulli occupation, $n_I\langle S_z^2\rangle - n_I^2 M_z^2$. The weight should therefore be
$v_z = (1-n_I)M_z^2 + (1-M_z)^2/3$. The printed value overestimates the longitudinal
fluctuations by a factor $1+n_I$. This factor is positive, so it changes the size of the
fluctuation term and never its direction: the direction is set by the sign alone. The corrected
curves below include both corrections.

%======================================================================
\section{Checks that do not use the derivation}\label{s:checks}

\subsection{Matrix SCBA}

The SCBA can be solved without the reduction to four scalars. The self-energy is iterated as a
$4\times4$ matrix, $\Sigma = \sum_t w_t V_t\langle G\rangle_{\mathrm{loc}}V_t$, where $V_t$ are the
exchange vertices with their second-cumulant weights and
$\langle G\rangle_{\mathrm{loc}} = N^{-2}\sum_{\bm k}[E+\I\eta-H_0(\bm k)-\langle V\rangle-\Sigma]^{-1}$, on the same
$k$ grid. The reduced equations with the sign of Eq.~(\ref{e:mok}) reproduce this matrix
solution to $10^{-7}$, for both couplings. With the sign of Eq.~(\ref{e:m}) they do not.

\subsection{Exact renormalised masses}\label{s:exact}

The disorder-averaged Green's function at $\bm k=0$ defines the exact self-energy, to all orders
in $J_I$:
\begin{equation}
\Sigma(E) = E - H_0(\bm 0) - \langle G(\bm 0, E)\rangle^{-1}.
\end{equation}
In each block, $H_0(\bm 0)=m\tau_z$ and $\langle G(\bm 0)\rangle$ is diagonal in the E/H basis, so
\begin{equation}
\bar m = m + \tfrac12\,\mathrm{Re}(\Sigma_{EE}-\Sigma_{HH}),\qquad
\bar\mu = -\tfrac12\,\mathrm{Re}(\Sigma_{EE}+\Sigma_{HH}),\qquad
\Gamma = -\tfrac12\,\mathrm{Im}(\Sigma_{EE}+\Sigma_{HH}).
\end{equation}
These include the mean field, like the SCBA's $\bar m_\sigma$ and $\bar\mu_\sigma$.
The broadening $\Gamma$ stays near zero when no state lies at $E$.

We computed $\langle G(\bm 0)\rangle_{aa} = \langle\phi_a|(E+\I\eta-H)^{-1}|\phi_a\rangle$:
\begin{itemize}
\item $\phi_a$ is the normalised $\bm k=0$ plane wave on orbital $a$;
\item $H$ is the Hamiltonian of a $160\times160$ periodic lattice with
  $\lfloor n_I\cdot160^2\rfloor$ impurities, and $\eta=0.006$;
\item the resolvent is expanded in Chebyshev polynomials of the rescaled Hamiltonian
  (Appendix~\ref{a:methods}).
\end{itemize}
Two independent samples agree to $10^{-3}$. Fig.~\ref{f:masses} and Table~\ref{t:masses} compare
the result with both versions of the SCBA.

\begin{figure}[t]
\centering
\fig{born_masses}
\caption{\label{f:masses}Renormalised masses: exact (circles), published SCBA (dashed),
corrected SCBA (solid), mean field (dotted). A block is topological for $\bar m<0$. Left:
same-$g$, $n_I=0.2$, $E_F=0$. Right: opposite-$g$, $n_I=0.4$, $E_F=0.02$. The exact masses come
from the disorder-averaged Green's function of a $160\times160$ periodic sample (no edges),
2 configurations, $\eta=0.006$, with no Born approximation. The SCBA curves use the grid of the
thesis program, $256\times64$ points of the Brillouin zone, with $\eta=10^{-3}$.}
\end{figure}

\begin{table}[t]
\centering
\caption{\label{t:masses}Renormalised masses, exact and in the two versions of the SCBA.}
\begin{tabular}{lcccc}
\toprule
case & $J_I$ & exact & published & corrected\\
\midrule
same-$g$, $\bar m_\up$, $E_F=0$ & 0.2 & $-0.1428$ & $-0.1312$ & $-0.1419$\\
 & 0.4 & $-0.1720$ & $-0.1121$ & $-0.1597$\\
 & 0.6 & $-0.2432$ & $-0.0705$ & $-0.1989$\\
same-$g$, $\bar m_\dn$, $E_F=0$ & 0.2 & $-0.1408$ & $-0.1319$ & $-0.1413$\\
 & 0.4 & $-0.1500$ & $-0.1175$ & $-0.1536$\\
 & 0.6 & $-0.1628$ & $-0.0960$ & $-0.1734$\\
opposite-$g$, $\bar m_\up$, $E_F=0.02$ & 0.30 & $-0.0302$ & $+0.0029$ & $-0.0313$\\
 & 0.35 & $-0.0147$ & $+0.0307$ & $-0.0164$\\
 & 0.40 & $+0.0002$ & $+0.0591$ & $-0.0021$\\
 & 0.45 & $+0.0145$ & $+0.0884$ & $+0.0098$\\
opposite-$g$, $\bar m_\dn$, $E_F=0.02$ & 0.40 & $-0.3377$ & $-0.2506$ & $-0.3422$\\
same-$g$, $\bar m_\dn$, $E_F=-0.05$ & 0.8 & $-0.1790$ & $-0.0605$ & $-0.2084$\\
 & 1.0 & $-0.1964$ & $-0.0176$ & $-0.2569$\\
 & 1.2 & $-0.2150$ & $+0.0331$ & $-0.3129$\\
 & 1.5 & $-0.2485$ & $+0.1235$ & $-0.3511$\\
\bottomrule
\end{tabular}
\end{table}

\begin{itemize}
\item \textbf{Direction.} Measured from the mean field, the exact masses move towards the
  topological side, like the corrected SCBA and unlike the published one.
\item \textbf{Weak disorder.} At same-$g$, $J_I=0.2$, the corrected SCBA is within $10^{-3}$ of
  the exact masses. The published one is off by $0.01$, in the opposite direction.
\item \textbf{The inversion.} At opposite-$g$, the exact spin-up mass changes sign at
  $J_I=0.40$, where the corrected SCBA puts it. The published SCBA puts it at $0.30$.
\item \textbf{Strong disorder.} The corrected SCBA remains close up to $J_I=0.6$. At same-$g$
  the two blocks drift apart from it in opposite directions. This is the third cumulant of the
  binary disorder, which the SCBA leaves out (below).
\item \textbf{The block that carries $G=1$.} At same-$g$ and $E_F=-0.05$, deep in the $G=1$
  plateau, the exact spin-down mass stays topological and its gap widens up to $J_I=1.5$. The
  broadening there stays below $0.005$: no state lies at $E_F$. The published SCBA inverts this
  block at $J_I\approx1.1$, which does not happen.
\item \textbf{Spin up at strong disorder.} Table~\ref{t:resonance} follows the line
  $E_F=-0.06$, through the onset of the $G=1$ plateau.
  \begin{itemize}
  \item For $J_I\gtrsim0.4$, $E_F$ lies in the spin-up band, and the spin-up states are strongly
    scattered.
  \item Near $J_I\approx0.87$, the $\bm k=0$ self-energy of that block goes through a resonance:
    the broadening peaks at $1.2$, and the mass diverges and changes sign.
  \item Beyond it, the block has a trivial gap around $E_F$: $\bar m_\up>0$, with a broadening of
    $0.01$ at $J_I=1.2$ and $0.006$ at $1.5$ ($E_F=-0.05$).
  \end{itemize}
  This is a strong-scattering effect, outside any Born approximation. The published SCBA
  reaches a trivial spin-up block through a smooth inversion, at $J_I=0.72$ on this line; the
  corrected one never does. Both predict $T_\up=0$ in this region, as observed. The spin-down
  block stays topological throughout, with no state at $E_F$.
\end{itemize}

\begin{table}[t]
\centering
\caption{\label{t:resonance}Exact $\bm k=0$ parameters along $E_F=-0.06$ (same-$g$, $n_I=0.2$),
across the onset of the $G=1$ plateau, with the published spin-up mass for comparison.}
\begin{tabular}{cccccc}
\toprule
$J_I$ & $\bar m_\up$ & $\Gamma_\up$ & published $\bar m_\up$ & $\bar m_\dn$ & $\Gamma_\dn$\\
\midrule
0.60 & $-0.270$ & 0.10 & $-0.055$ & $-0.164$ & 0.0004\\
0.70 & $-0.273$ & 0.21 & $-0.010$ & $-0.171$ & 0.0005\\
0.80 & $-0.569$ & 0.37 & $+0.053$ & $-0.179$ & 0.0006\\
0.85 & $-1.212$ & 1.24 & $+0.088$ & $-0.183$ & 0.0006\\
0.90 & $+1.176$ & 0.38 & $+0.113$ & $-0.188$ & 0.0007\\
0.95 & $+0.552$ & 0.08 & $+0.124$ & $-0.192$ & 0.0008\\
1.00 & $+0.369$ & 0.04 & $+0.132$ & $-0.197$ & 0.0008\\
\bottomrule
\end{tabular}
\end{table}

\paragraph{Convergence.} Fig.~\ref{f:masses256} repeats Fig.~\ref{f:masses} with a
$256\times256$ periodic sample, and with the SCBA on a $256\times256$ grid. Nothing moves:
\begin{itemize}
\item the exact masses change by at most $2\times10^{-3}$ (at $J_I=0.6$);
\item the SCBA changes by at most $6\times10^{-4}$, and by at most $2\times10^{-3}$ when $\eta$ goes
  from $10^{-3}$ to $6\times10^{-3}$;
\item halving $\eta$ to $0.003$ changes the exact masses by at most $2\times10^{-3}$.
\end{itemize}
The remaining same-$g$ offset is therefore not numerical.

\begin{figure}[t]
\centering
\fig{born_masses_L256_k256}
\caption{\label{f:masses256}Fig.~\ref{f:masses} with a $256\times256$ periodic sample for the
exact masses (2 configurations, $\eta=0.006$) and a $256\times256$ grid for the SCBA.}
\end{figure}

\paragraph{Where the same-$g$ offset comes from.} The SCBA keeps only the first two moments of
the disorder (Step~2 of Sec.~\ref{s:derivation}). The impurity potential, present or absent on
each site, also has a third cumulant,
\begin{equation}
\big\langle(\delta U)^3\big\rangle=\pm\,n_I(1-n_I)(1-2n_I)\,J_I^3 ,
\end{equation}
which enters the self-energy at third order.
\begin{itemize}
\item \textbf{Gaussian test.} We replaced the impurities by a Gaussian potential on every site,
  with the same mean and variance, so all higher cumulants vanish. The exact masses then fall
  onto the corrected SCBA, within $3\times10^{-3}$ up to $J_I=0.5$ (Table~\ref{t:gauss}).
\item \textbf{Opposite signs for the two spins.} At same-$g$, spin up sees $+J_I$ and spin down
  $-J_I$, so the third-order term changes sign with the spin. That is why the two blocks leave
  the SCBA in opposite directions.
\item \textbf{Why opposite-$g$ is so close.} Relative to the variance, the third cumulant scales
  as $(1-2n_I)J_I$: $0.6\,J_I$ at $n_I=0.2$, but $0.2\,J_I$ at $n_I=0.4$.
\end{itemize}

\begin{table}[t]
\centering
\caption{\label{t:gauss}Same-$g$, $n_I=0.2$, $E_F=0$: exact masses with the binary impurities
and with a Gaussian potential of the same mean and variance ($160\times160$ periodic sample,
2 configurations, $\eta=0.006$), against the corrected SCBA.}
\begin{tabular}{ccccccc}
\toprule
 & \multicolumn{3}{c}{$\bar m_\up$} & \multicolumn{3}{c}{$\bar m_\dn$}\\
$J_I$ & binary & Gaussian & SCBA & binary & Gaussian & SCBA\\
\midrule
0.3 & $-0.1525$ & $-0.1491$ & $-0.1487$ & $-0.1448$ & $-0.1466$ & $-0.1465$\\
0.4 & $-0.1720$ & $-0.1628$ & $-0.1597$ & $-0.1500$ & $-0.1541$ & $-0.1536$\\
0.5 & $-0.2082$ & $-0.1777$ & $-0.1795$ & $-0.1560$ & $-0.1635$ & $-0.1625$\\
0.6 & $-0.2432$ & $-0.1899$ & $-0.1989$ & $-0.1628$ & $-0.1759$ & $-0.1734$\\
\bottomrule
\end{tabular}
\end{table}

\subsection{Spin-resolved transport}\label{s:spin}

With $M_z=1$, the transmission splits exactly into $T=T_\up+T_\dn$. We computed it on the ribbons
of the paper's maps: $400\times64$ sites, $n_I\cdot256\cdot64$ impurities in the central 256
columns, 8 configurations per point. Table~\ref{t:spin} lists the points where the two versions
disagree on the topology of a block.

\begin{table}[t]
\centering
\caption{\label{t:spin}Exact spin-resolved transmissions and the renormalised mass of the relevant
block in the two versions of the SCBA. Means over 8 configurations ($W=64$, $128$; ribbons of
$400$ sites) or 3 configurations ($W=256$; disordered length $256$).}
\begin{tabular}{lcccccc}
\toprule
case & $J_I$ & $E_F$ & $T_\up$ & $T_\dn$ & published & corrected\\
\midrule
same-$g$, $W=64$ & 1.2 & $-0.05$ & 0.000 & 1.000 & $\bar m_\dn=+0.033$ (trivial) & $\bar m_\dn=-0.313$\\
same-$g$, $W=64$ & 1.3 & $-0.12$ & 0.000 & 1.000 & $\bar m_\dn=+0.080$ (trivial) & $\bar m_\dn=-0.314$\\
opposite-$g$, $W=64$ & 0.30 & 0.02 & 0.840 & 1.000 & $\bar m_\up=+0.003$ & $\bar m_\up=-0.031$\\
opposite-$g$, $W=64$ & 0.35 & 0.02 & 0.005 & 1.000 & $\bar m_\up=+0.031$ (trivial) & $\bar m_\up=-0.016$\\
opposite-$g$, $W=128$ & 0.30 & 0.02 & 0.997 & 1.000 & & \\
opposite-$g$, $W=128$ & 0.35 & 0.02 & 0.562 & 1.000 & & \\
opposite-$g$, $W=128$ & 0.40 & 0.02 & 0.012 & 1.000 & $\bar m_\up=+0.059$ & $\bar m_\up=-0.002$\\
opposite-$g$, $W=256$ & 0.35 & 0.02 & 0.989 & 1.000 & $\bar m_\up=+0.031$ (trivial) & $\bar m_\up=-0.016$\\
opposite-$g$, $W=256$ & 0.40 & 0.02 & 0.339 & 1.000 & $\bar m_\up=+0.059$ & $\bar m_\up=-0.002$\\
\bottomrule
\end{tabular}
\end{table}

\begin{itemize}
\item \textbf{Same-$g$.} The spin-down block carries a quantized channel, in 8 configurations out
  of 8, where the published masses make it trivial. A trivial block cannot carry a quantized
  chiral channel.
\item \textbf{Opposite-$g$.} The spin-up channel vanishes earlier than the exact inversion, and
  later when the ribbon is wider. A trivial block would carry no channel at any width. This
  finite-width effect is the subject of Sec.~\ref{s:width}.
\end{itemize}

%======================================================================
\section{Impact on the Born phase diagrams}\label{s:impact}

Fig.~\ref{f:phase} draws both versions of the Born lines over the published conductance maps.
These are the data of Fig.~1 of Ref.~\cite{g2g1}, averaged over 101 configurations
(opposite-$g$) and 50 configurations (same-$g$), in the colours of the thesis. Table~\ref{t:phase}
counts the points of the $G=1$ plateau that each version explains. A block counts for one
channel if it is topological and $E_F$ lies in its gap (``localised'' reading). Alternatively it
counts if it is topological, whatever $E_F$ (``topological'' reading).

\begin{figure}[p]
\centering
\fig{born_phase}
\caption{\label{f:phase}Born lines over the published conductance maps (Fig.~1 of
Ref.~\cite{g2g1}): published (left) and corrected SCBA (right). Light blue: $|\langle G\rangle-1|<0.05$; blue:
$|\langle G\rangle-2|<0.05$. Solid lines: spin-up band edges; dashed: spin-down band edges; dotted:
$\bar m_\sigma=0$.}
\end{figure}

\begin{table}[t]
\centering
\caption{\label{t:phase}Share of the $G=1$ plateau ($|\langle G\rangle-1|<0.05$) of the published maps
predicted at $G=1$ by the Born lines.}
\begin{tabular}{lccccc}
\toprule
 & & \multicolumn{2}{c}{localised reading} & \multicolumn{2}{c}{topological reading}\\
map & $G=1$ points & published & corrected & published & corrected\\
\midrule
same-$g$, $n_I=0.2$ & 1065 & 20\,\% & 99.9\,\% & 39\,\% & 0\,\%\\
opposite-$g$, $n_I=0.4$ & 122 & 94\,\% & 88\,\% & 96\,\% & 59\,\%\\
\bottomrule
\end{tabular}
\end{table}

\needspace{8\baselineskip}
\paragraph{Same-$g$.}
\begin{itemize}
\item \textbf{Published lines.} Both masses change sign: $\bar m_\up=0$ at $J_I\approx0.65$--$0.95$,
  then $\bar m_\dn=0$ at $J_I\approx1.0$--$1.2$. Beyond the second line no block is topological,
  so $G$ should vanish. The published map shows $G=1$ there, up to $J_I=1.5$: 566 of its 1065
  $G=1$ points lie where the published masses make both blocks trivial. The shaded $G=1$ region
  of the paper's Fig.~4 covers only a small part of the plateau.
\item \textbf{Corrected lines.} No mass changes sign. They locate where the spin-up channel
  starts to fall, at the spin-up band edge, and keep spin down topological, as the exact masses
  do. The strong-disorder end of the spin-up behaviour (Sec.~\ref{s:exact}) is beyond them.
\end{itemize}
Table~\ref{t:onset} follows the published map along lines of constant $E_F$:
\begin{itemize}
\item $G$ leaves 2 at the corrected spin-up band edge.
\item The plateau begins 0.2--0.4 further, along a line with the same slope.
\item The published $\bar m_\up=0$ line is almost vertical ($J_I=0.65$--$0.83$). It lies within
  0.04 of the onset of the plateau for $-0.11\lesssim E_F\lesssim-0.06$, and 0.08--0.2 away
  elsewhere.
\item On the line $E_F=-0.06$, the exact spin-up block does turn trivial, at $J_I\approx0.9$,
  through the resonance of Sec.~\ref{s:exact}. The paper's picture of a trivial spin-up block
  therefore holds at strong disorder, but not through the mechanism, nor at the place, that the
  published line gives.
\end{itemize}

\begin{table}[t]
\centering
\caption{\label{t:onset}Same-$g$ published map along lines of constant $E_F$: where $\langle G\rangle$ leaves 2
($<1.95$) and where the $G=1$ plateau begins ($|\langle G\rangle-1|<0.05$), against the corrected spin-up
band edge and the published $\bar m_\up=0$ line.}
\begin{tabular}{ccccc}
\toprule
$E_F$ & $G$ leaves 2 & plateau begins & corrected E$\up$ edge & published $\bar m_\up=0$\\
\midrule
$-0.109$ & 0.24 & 0.62 & 0.14 & 0.65\\
$-0.084$ & 0.31 & 0.64 & 0.26 & 0.68\\
$-0.058$ & 0.38 & 0.76 & 0.38 & 0.72\\
$-0.033$ & 0.48 & 0.83 & 0.48 & 0.75\\
$-0.008$ & 0.52 & 0.93 & 0.59 & 0.78\\
$+0.018$ & 0.57 & 0.98 & 0.70 & 0.81\\
\bottomrule
\end{tabular}
\end{table}

\paragraph{Opposite-$g$.}
\begin{itemize}
\item \textbf{Published line.} $\bar m_\up=0$ at $J_I\approx0.29$ lies at the left edge of the
  $G=1$ region, as the paper describes.
\item \textbf{Corrected line.} It lies at $J_I\approx0.41$--$0.45$, inside that region.
\end{itemize}
Sec.~\ref{s:width} shows that the left edge of the $G=1$ region is not the inversion but a
finite-width effect.

%======================================================================
\section{Why the published inversion line matched the antiferromagnetic map}\label{s:width}

\subsection{Overlap of the edge states}

An edge state of a BHZ block with mass $m$ decays into the bulk over $\xi=1/\lambda_-$. Here
$\lambda_\pm^2$ are the roots of $(E_0-D\lambda^2)^2=(m+B\lambda^2)^2-A^2\lambda^2$ at the edge-state energy
$E_0=-mD/B$~\cite{zhou2008}. For $|m|\ll A^2|B|/(B^2-D^2)$,
\begin{equation}
\xi \simeq \frac{A|B|}{|m|\sqrt{B^2-D^2}} \simeq \frac{1.50}{|m|}\ \text{sites}.
\end{equation}
This gives $\xi=10.7$ sites for the clean mass, 50 sites for $|m|=0.03$ and 100 sites for
$|m|=0.015$. In a ribbon of width $W$, the two counter-propagating spin-up edge states overlap
once $\xi$ becomes comparable with $W$. The disorder then backscatters between them across the
ribbon, and the block's channel is lost while the block is still topological~\cite{zhou2008}. For
the paper's $W=64$, this happens at $|\bar m_\up|\approx0.015$--$0.03$ ($\xi\approx50$--$100$
sites), i.e. $J_I\approx0.30$--$0.35$ at $n_I=0.4$ (Table~\ref{t:masses}). That is where the published map leaves $G=2$, where the
published $\bar m_\up=0$ line happens to lie.

\subsection{Wider ribbons}

Fig.~\ref{f:width} (left) follows both blocks along $E_F=0.02$ for $W=64$, $128$ and $256$.
The disordered length is 256, with 8, 3 and 3 configurations.
\begin{itemize}
\item \textbf{The spin-up channel disappears later in wider ribbons.} $T_\up=1/2$ at
  $J_I\approx0.32$ ($W=64$), $0.355$ ($W=128$) and $0.39$ ($W=256$). The transition converges on
  the exact inversion, $J_I=0.40$ (Table~\ref{t:masses}), and the channel is lost when $\xi$
  becomes comparable with $W$.
\item \textbf{The published inversion is ruled out.} At $W=256$ and $J_I=0.35$, the spin-up block
  carries $T_\up=0.989$, the same within $0.002$ in the three configurations. The published SCBA
  makes this block trivial there, with a gap of $0.06$. A trivial block carries no channel at
  any width.
\item \textbf{The spin-down channel also survives longer.} At $W=256$, $T_\dn=1.000$ up to
  $J_I=0.475$, and it stays above $0.8$ on average up to $0.6$. At $W=64$ it is already $0.89$
  at $J_I=0.45$ and $0.60$ at $0.5$. Its loss in narrow ribbons comes from the impurity states
  that fill its gap (Sec.~\ref{s:origin}), which connect the two edges more easily when they
  are close.
\item \textbf{The opposite-$g$ $G=1$ window moves and widens.} At $W=256$ it starts at the exact
  inversion, $J_I\approx0.40$, and extends beyond $J_I=0.6$.
\end{itemize}
Fig.~\ref{f:width} (right) follows the same-$g$ line $E_F=-0.06$.
\begin{itemize}
\item The spin-down block carries $T_\dn=1.000$ at every $J_I$ up to 1.5 and at every width, far
  beyond the published $\bar m_\dn=0$ line.
\item $T_\up$ starts to fall at the corrected spin-up band edge ($J_I\approx0.38$) at every width.
  It then crosses a diffusive region, where its value grows with $W$, and vanishes from
  $J_I\approx0.8$ at every width.
\end{itemize}
Fig.~\ref{f:widephase} shows the full maps. The caroli maps at $W=64$ reproduce the published
ones. At $W=128$:
\begin{itemize}
\item the same-$g$ plateau is unchanged;
\item the opposite-$g$ $G=2$ region extends to the corrected inversion, and the $G=1$ region
  sits in the wedge beyond it, where the spin-up gap is trivial.
\end{itemize}

\begin{figure}[t]
\centering
\fig{born_wide_phase}
\caption{\label{f:widephase}$\langle G\rangle$ maps in the colours of the thesis: published map
($W=64$, Fortran, 101 and 50 configurations), and caroli on the same disordered length for
$W=64$ (8 configurations) and $W=128$ (3 configurations). Lines: corrected SCBA.}
\end{figure}

\begin{figure}[t]
\centering
\fig{born_width}
\caption{\label{f:width}Transmissions of the two blocks along a line of constant $E_F$ for
$W=64$, $128$ and $256$ (8, 3 and 3 configurations; disordered length 256). Left: opposite-$g$,
$E_F=0.02$, with the published and corrected spin-up inversions. Right: same-$g$, $E_F=-0.06$,
with the corrected spin-up band edge and the published inversions of both blocks.}
\end{figure}

\subsection{The coincidence}

The wrong sign moves the Born inversion from $J_I=0.40$ to $0.30$. The finite width moves the
observed loss of the spin-up channel from $0.40$ to about $0.32$--$0.35$ at $W=64$. The two
shifts have the same direction and similar sizes, so the published line happened to fall on the
edge of the published map. The same wrong sign predicts that both blocks become trivial at
same-$g$, which the published map itself contradicts (Sec.~\ref{s:impact}).

%======================================================================
\section{Origin and stability of the $G=1$ phase}\label{s:origin}

Fig.~\ref{f:spinmaps} resolves the conductance by spin. In each block, the corrected band edges
of that block bound the region where it carries one channel.

\begin{figure}[t]
\centering
\fig{born_spin_maps}
\caption{\label{f:spinmaps}Spin-resolved transmissions $\langle T_\up\rangle$ (left) and $\langle T_\dn\rangle$
(right), $W=128$, 3 configurations, with the corrected band edges of each block, in the
colours of Fig.~\ref{f:phase}. Light blue: the block carries one channel
($|\langle T_\sigma\rangle-1|<0.05$); black: none. Top: opposite-$g$, $n_I=0.4$. Bottom:
same-$g$, $n_I=0.2$.}
\end{figure}

\paragraph{Two independent Chern insulators.} With $M_z=1$, $s_z$ is conserved and each block is a
two-dimensional system without time-reversal symmetry (class A). Its two-terminal conductance is
quantized to the Chern number of the occupied states whenever the states at $E_F$ are
localised. A block therefore contributes one channel as long as $E_F$ lies above the delocalised
states of its valence band and below those of its conduction band.

\paragraph{Same-$g$: how the phase appears.}
\begin{itemize}
\item The mean exchange field acts as a spin-dependent chemical potential: the spin-up bands
  rise by $n_I J_I$ and the spin-down bands sink by $n_I J_I$.
\item The fluctuations widen both gaps (Fig.~\ref{f:masses}).
\item Near $E_F=0$, the Fermi level enters the spin-up valence band as $J_I$ grows: $T_\up$ starts
  to fall at the corrected spin-up band edge.
\item Once $E_F$ lies below the delocalised spin-up states, the spin-up states at $E_F$ are
  localised and carry no Chern number, so $T_\up=0$. The grey crossover of the map is the
  passage through those states.
\item At stronger disorder ($J_I\gtrsim0.9$ on $E_F=-0.06$), the spin-up block goes through a
  resonance and becomes a trivial insulator at $E_F$ (Table~\ref{t:resonance}). $T_\up$ stays
  zero.
\item The spin-down gap still contains $E_F$, so $T_\dn=1$.
\item The plateau ends only when $E_F$ reaches the spin-down conduction band, at the top right
  of the map.
\end{itemize}
On the plateau of the $W=64$ maps (8 configurations), $T_\dn$ is within $10^{-3}$ of 1 in $97\,\%$
of the single configurations, and $T_\up<10^{-3}$ in $92\,\%$.

\paragraph{Why it is so stable.}
\begin{itemize}
\item The $G=1$ phase is a Chern insulator of total Chern number 1. Each edge carries a single
  chiral mode. There is no counter-propagating mode to scatter into, in either spin, so no
  disorder can backscatter it as long as the bulk states at $E_F$ are localised.
\item The $G=2$ phase is different: each edge carries two counter-propagating modes, one per
  spin, and only $s_z$ conservation keeps them apart.
\item Tilting the moments shows the difference. With $M_z=0.8$ the in-plane components couple
  the spins. On the thesis map at $M_z=0.8$ (same-$g$, $n_I=0.2$, 56 configurations), the
  $G=2$ plateau shrinks from $23\,\%$ to $3.8\,\%$ of the map, while the $G=1$ plateau keeps
  $13.9\,\%$ (against $26\,\%$ at $M_z=1$).
\end{itemize}

\paragraph{Opposite-$g$.}
\begin{itemize}
\item The mean field is a spin-dependent mass, so the spin-up block does change sign, at
  $J_I\simeq0.40$ for $n_I=0.4$. The fluctuations delay this inversion relative to the mean
  field ($0.34$).
\item In Fig.~\ref{f:spinmaps}, $T_\up=1$ fills the corrected topological gap of the spin-up
  block: the wedge between its two band edges, which closes at $J_I\approx0.41$,
  $E_F\approx0.02$. $T_\up=0$ fills the trivial gap that opens beyond.
\item The $G=1$ phase is the region where the spin-up gap is trivial while spin down still
  conducts.
\item The spin-down block becomes more topological, but at $n_I=0.4$ its gap fills with
  impurity states for $J_I\gtrsim0.45$. The exact broadening at $E_F=0.02$ rises from $0.001$ at
  $J_I=0.35$ to $0.013$, $0.024$ and $0.048$ at $J_I=0.45$, $0.5$ and $0.6$. At $W=64$, $T_\dn$
  drops to $0.89$, $0.60$ and $0.32$.
\item In the paper's 64-site ribbons, the $G=1$ window is therefore narrow at this
  concentration. It widens in wider ribbons, where both channels survive longer
  (Sec.~\ref{s:width}).
\end{itemize}

%======================================================================
\section{What changes for the paper}\label{s:addendum}

\paragraph{Unchanged.}
\begin{itemize}
\item The numerical results: conductance maps (Fig.~1), local currents (Fig.~2) and local
  densities of states (Fig.~3).
\item The main conclusion: a disorder-induced $G=1$ Chern-insulator phase for both couplings,
  carried by the spin-down edge state and caused by the suppression of one spin channel through
  selective scattering and localisation.
\end{itemize}

\paragraph{Changed.}
\begin{itemize}
\item The sign of the mass equations of the SCBA, and the weight $v_z$.
\item Fig.~4 and the paragraphs that interpret it:
  \begin{itemize}
  \item same-$g$: the spin-down mass never changes sign, whereas the published lines invert
    it at $J_I\approx1.0$--$1.2$. The plateau begins where the spin-up states at $E_F$ stop
    conducting, after $E_F$ has entered the spin-up band, shifted by the exchange field. The
    spin-up block becomes trivial only at stronger disorder, through a resonance beyond the
    Born approximation;
  \item opposite-$g$: the spin-up inversion lies at $J_I\simeq0.40$ for $n_I=0.4$, and the left
    edge of the $G=1$ region in the maps is set by the ribbon width.
  \end{itemize}
\end{itemize}

\paragraph{Draft addendum.}
\begin{quote}\small
In the self-consistent Born analysis of Sec.~IV\,B, the definition of the renormalised
parameters should read
$\mathrm{Re}\,\Sigma = \mathrm{diag}[(\bar m_\up-m_\up)\tau_z-(\bar\mu_\up-\mu_\up)\tau_0,\;(\bar m_\dn-m_\dn)\tau_z-(\bar\mu_\dn-\mu_\dn)\tau_0]$,
so that the mass equations carry a plus sign in front of the sum. In addition, the
longitudinal weight is $v_z=(1-n_I)M_z^2+(1-M_z)^2/3$. The corrected Born phase diagrams are
shown below.

With the correct sign, the fluctuations renormalise both masses towards the topological side,
as in the topological Anderson insulator, and the same-$g$ spin-down mass does not change sign. The
$G=1$ phase of the same-$g$ case arises because the spin-dependent shift of the bands brings
the Fermi level into the spin-up band, whose states are localised by the disorder. At stronger
disorder, the spin-up block becomes trivial through resonant scattering, beyond the Born
approximation. Meanwhile the spin-down block keeps a topological gap and its chiral edge state,
as shown in Fig.~2. In
the opposite-$g$ case, the spin-up inversion occurs at $J_I\approx0.40$ for $n_I=0.4$.

The numerical results and the conclusions of the paper are unchanged.
\end{quote}

%======================================================================
\appendix
\section{Methods and reproducibility}\label{a:methods}

All calculations use the Rust toolbox \texttt{caroli}, which reproduces the thesis programs on
their own test cases. Each item below is an example, run with
\begin{verbatim}
  cargo run --release -p caroli-validation --example <name>
\end{verbatim}
\begin{itemize}
\item \texttt{born\_lines}: published and corrected Born parameters, on the grids of the
  published maps.
\item \texttt{born\_self\_energy}: exact renormalised masses (Sec.~\ref{s:exact}).
\item \texttt{born\_spin\_channels}: spin-resolved transmissions of Table~\ref{t:spin}.
\item \texttt{born\_wide\_maps}: spin-resolved maps and cuts for $W=64,128,256$.
\end{itemize}
The figures come from the plotting scripts \texttt{plot\_born\_phase.py},
\texttt{plot\_born\_masses.py} and \texttt{plot\_born\_wide.py} in the same folder.

\paragraph{Chebyshev expansion of the resolvent.} With $\tilde H=(H-c)/h$ rescaled into $(-1,1)$,
$\tilde z = (E+\I\eta-c)/h$, $s=\sqrt{\tilde z^2-1}$ and $t=\tilde z-s$ chosen with $|t|<1$,
\begin{equation}
\langle\phi|(E+\I\eta-H)^{-1}|\phi\rangle = \frac{1}{h\,s}\Big[\mu_0+2\sum_{n\ge1}t^n\mu_n\Big],\qquad
\mu_n = \langle\phi|T_n(\tilde H)|\phi\rangle .
\end{equation}
The moments are obtained two per product with $\tilde H$, from
$\mu_{2n}=2\langle T_n|T_n\rangle-\mu_0$ and $\mu_{2n+1}=2\langle T_n|T_{n+1}\rangle-\mu_1$. The series is
truncated where $|t|^n<10^{-10}$.

\paragraph{Spin-resolved transport.} With $M_z=1$ each block is a two-orbital BHZ model with its
own disorder. Its transmission is computed by recursive Green's functions, with the lead
self-energies from the propagating and evanescent modes. The blocks reproduce the four-orbital
calculation to $10^{-12}$.

\begin{thebibliography}{9}
\bibitem{g2g1} L.~Raymond, A.~D.~Verga, and A.~Demion, \emph{Anomalous quantum Hall effect induced
  by disorder in topological insulators}, \todo{journal reference}.
\bibitem{thesis} A.~Demion, \emph{Transport électronique dans le graphène et les isolants
  topologiques 2D en présence de désordre magnétique}, PhD thesis, Aix-Marseille Université
  (2015).
\bibitem{li2009} J.~Li, R.-L.~Chu, J.~K.~Jain, and S.-Q.~Shen, Phys.\ Rev.\ Lett.\ \textbf{102},
  136806 (2009).
\bibitem{groth2009} C.~W.~Groth, M.~Wimmer, A.~R.~Akhmerov, J.~Tworzyd\l{}o, and
  C.~W.~J.~Beenakker, Phys.\ Rev.\ Lett.\ \textbf{103}, 196805 (2009).
\bibitem{jiang2009} H.~Jiang, L.~Wang, Q.-F.~Sun, and X.~C.~Xie, Phys.\ Rev.\ B \textbf{80},
  165316 (2009).
\bibitem{zhou2008} B.~Zhou, H.-Z.~Lu, R.-L.~Chu, S.-Q.~Shen, and Q.~Niu, Phys.\ Rev.\ Lett.\
  \textbf{101}, 246807 (2008).
\end{thebibliography}

\end{document}
